\chapter{Introduction}
\label{ch:intro}

\section{Background and motivation}

Every product we use -- a packet of biscuits, a litre of milk, a steel beam -- carries an
environmental burden that is spread across a long chain of activities: growing or extracting raw
materials, processing, packaging, transport, retail, use and disposal. \emph{Life Cycle
Assessment} (LCA), standardised in ISO~14040 and ISO~14044 \citep{iso14040,iso14044}, is the
accepted method for adding up these burdens in a consistent way. LCA results now drive
eco-labels, Environmental Product Declarations (EPDs), corporate carbon reporting, green public
procurement and product redesign \citep{hellweg2014,finnveden2009}.

The food system alone is responsible for roughly one third of global anthropogenic greenhouse-gas
emissions \citep{crippa2021}, and the footprints of individual foods differ by orders of magnitude
-- beef can emit more than 50 times as much per kilogram as many plant foods \citep{poore2018}.
Choosing the right products, recipes and suppliers is therefore one of the largest levers
available to manufacturers, retailers and consumers. Yet acting on this lever requires
footprint numbers for very large numbers of products.

\section{The problem: LCA is accurate but slow}

A full process-based LCA requires a detailed \emph{life cycle inventory} (LCI): every material,
energy and emission flow for every process in the supply chain. Assembling this inventory is the
most time-consuming and expensive step \citep{reap2008}. Commercial background databases such as
ecoinvent \citep{wernet2016} are licensed, and even with database access a practitioner must
(i) map every item of a bill of materials or recipe to the right database process, and
(ii) fill gaps where no matching data exist. A single study can take weeks to months.

This creates two practical bottlenecks that this thesis targets:

\begin{enumerate}[label=\textbf{B\arabic*.}]
  \item \textbf{The estimation bottleneck.} Early in design, or when screening a portfolio of
  thousands of products, decision makers need an \emph{approximate but reliable} footprint long
  before a full inventory exists. The common shortcut -- borrowing the average of a product
  category -- is crude and gives no indication of its own uncertainty.
  \item \textbf{The mapping bottleneck.} Translating free-text descriptions (``frozen breaded
  chicken nuggets'', ``pâte feuilletée pur beurre'') into the correct database entries is
  manual, repetitive, language-dependent and error-prone.
\end{enumerate}

Artificial Intelligence and Machine Learning (AI/ML) are natural candidates for both: supervised
learning can learn the mapping from easily available product attributes to impacts, and modern
natural-language-processing (NLP) models can measure the semantic similarity of product
descriptions across languages. The question is \emph{how} to integrate them so that the result
is accurate, honest about its uncertainty, interpretable to an LCA practitioner and usable in
practice.

\section{Aim and objectives}

\textbf{Aim.} To design, implement and critically evaluate an AI/ML framework that delivers fast,
explainable and uncertainty-aware LCA screening estimates and automates the mapping of product
descriptions to LCA datasets, using only openly available data.

\noindent\textbf{Objectives.}
\begin{enumerate}[label=O\arabic*.]
  \item Review how AI/ML has been applied across the LCA phases and identify the gap addressed.
  \item Build a clean, leakage-free modelling dataset from the open AGRIBALYSE~3.1 database.
  \item Develop a surrogate impact model (Module~A) and benchmark it against practitioner
        baselines under both optimistic and realistic (unseen-category) validation.
  \item Quantify uncertainty with distribution-free conformal prediction intervals.
  \item Explain the model's predictions with SHAP and check them against LCA domain knowledge.
  \item Develop and benchmark a multilingual semantic mapping engine (Module~B).
  \item Integrate both modules into an interactive screening tool with an applicability-domain check.
\end{enumerate}

\section{Research questions}

\begin{description}[style=nextline,leftmargin=1.2cm]
  \item[RQ1] How accurately can ML predict product-level LCA impacts from pre-LCA information,
  compared with the category-average proxy used in practice?
  \item[RQ2] How much does accuracy degrade for product categories absent from training, and can
  the model's uncertainty be reported honestly?
  \item[RQ3] Do the model's explanations agree with established LCA knowledge?
  \item[RQ4] How well can sentence-embedding models map free-text product descriptions to LCA
  datasets, including across languages?
\end{description}

\section{Contributions}

\begin{itemize}
  \item A reproducible, open-data pipeline (AGRIBALYSE~3.1) with an explicit \emph{no-leakage}
        rule: no stage-wise or ingredient-wise LCA result is ever used as a model input.
  \item A two-scheme evaluation (random vs.\ leave-sub-group-out cross-validation) that separates
        interpolation from genuine generalisation -- a distinction often missing in ML-for-LCA work.
  \item Calibrated 90\% prediction intervals via Mondrian conformal prediction, evaluated in a
        nested design.
  \item A five-encoder benchmark for LCA dataset mapping in two realistic tasks (cross-lingual and
        terse queries).
  \item A working screening tool (single product and recipe/bill-of-materials modes) with
        out-of-domain warnings.
  \item An honest account of negative results (e.g.\ embedding-based stacking did not help) and
        limitations.
\end{itemize}

\section{Scope}

The case study is food products, chosen because (a) a high-quality, openly licensed, product-level
LCA database exists, (b) food is a major source of emissions, and (c) the product space is large
and diverse. The framework itself is domain-agnostic: any LCA database with product descriptors
and impact results can replace AGRIBALYSE. The tool is intended for \emph{screening and
prioritisation}, not for ISO-compliant comparative assertions or EPDs.

\section{Organisation of the thesis}

Chapter~\ref{ch:lit} reviews LCA fundamentals and the literature on AI/ML in LCA.
Chapter~\ref{ch:data} describes the data and exploratory analysis. Chapter~\ref{ch:method}
presents the methodology of both modules. Chapter~\ref{ch:results} reports the experiments.
Chapter~\ref{ch:tool} describes the prototype tool. Chapter~\ref{ch:discussion} discusses
implications, threats to validity and limitations, and Chapter~\ref{ch:conclusion} concludes with
future work.
